% ============================================================
%  part1/ch01-map.tex
%  Chapter 1: The Map Is Not the Territory
% ============================================================

\chapter{The Map Is Not the Territory}
\label{ch:map}

\epigraph{%
  The map is not the territory,\\
  and the name is not the thing named.%
}{Alfred Korzybski}

\bigskip

\begin{WhyMatters}
  The map example is not about cartography.
  It is the first appearance of a central theme:
  \emph{useful representations depend on forgetting.}
  Later chapters will show that detector measurements, memories,
  and cosmological observations obey exactly the same principle.
  The mathematics introduced here---projection, kernel,
  quotient---will reappear in every subsequent part of the book.
\end{WhyMatters}

\section{A useful object that is always wrong}

A map is a strange object.

It is useful precisely because it is wrong.

A road map does not show trees, clouds, buildings, birds, or people.
A topographic map ignores political borders.
A political map ignores elevation.
A transit map distorts distance while preserving connectivity.

Every successful map leaves something out.

At first glance this seems obvious.
No map could contain every detail of the territory it represents.
A perfect map would be as large and as complicated as the territory
itself and would therefore fail in its purpose.

But this observation has a deeper consequence.

The usefulness of a map depends not merely on what it contains,
but on what it discards.

A traveler cares about roads.
A geologist cares about elevation.
A city planner cares about infrastructure.

The same territory gives rise to many different maps because
different tasks require different projections.

\section{The projection}

\begin{TechNote}[Projection]
  Throughout this book, a \emph{projection} need not be linear.
  The term is used in the broader sense of any map
  \[
    \pi : X \to Y
  \]
  that preserves some structure while discarding other structure.
  Quotient maps, detector outputs, coarse-grainings, and
  observational restrictions are all projections in this sense.
\end{TechNote}

Let $X$ denote a territory and let $M$ denote a map.
The map is a projection
\[
  \pi : X \to M.
\]
The projection preserves some structure and discards the rest.

Suppose two distinct features of the territory are represented
identically on the map:
\[
  \pi(x_1) = \pi(x_2), \qquad x_1 \neq x_2.
\]
The map cannot distinguish between them.
From the perspective of the map, these two features are equivalent.

The territory has therefore been partitioned into equivalence classes:
\[
  x_1 \sim_\pi x_2
  \iff
  \pi(x_1) = \pi(x_2).
\]

\begin{definition}[Projection equivalence]
  \label{def:proj-equiv}
  Given a map $\pi : X \to M$, two points $x_1, x_2 \in X$ are
  \emph{projection-equivalent} if $\pi(x_1) = \pi(x_2)$.
  The equivalence class of $x$ is denoted $[x]_\pi$.
\end{definition}

The map is then a faithful representation of the quotient space:
\[
  M \cong X / {\sim_\pi}.
\]

\begin{ForwardRef}
  The equivalence relation $x_1 \sim_\pi x_2$ introduced here
  will later reappear as \emph{observational equivalence} in
  detector theory (Chapter~\ref{ch:detector}) and as
  \emph{horizon equivalence} in cosmology
  (Chapter~\ref{ch:horizons} and Part~VII).
  In all three cases the same mathematics governs what an
  observer cannot distinguish.
\end{ForwardRef}

\section{The kernel}

The information lost by a projection can be measured precisely.

\begin{definition}[Kernel of a projection]
  \label{def:kernel}
  Let $\pi : X \to M$ be a projection.
  The \emph{kernel} of $\pi$ is the collection of all pairs that
  $\pi$ fails to distinguish:
  \[
    \ker(\pi)
    = \{(x_1, x_2) \in X \times X : \pi(x_1) = \pi(x_2),\;
      x_1 \neq x_2\}.
  \]
  A projection is \emph{injective} (loses no information) if and
  only if $\ker(\pi) = \emptyset$.
\end{definition}

Every useful map has a nontrivial kernel.

Two cities at the same elevation appear identically on a topographic
contour.
Two neighborhoods in the same postal district appear identically on
a postal map.
Two points equidistant from the center appear identically on a
radial transit diagram.

The kernel is not a defect.
It is the mechanism by which the map achieves its purpose.

\begin{Misconception}
  Information loss does not imply ignorance.
  A map loses information yet remains useful.
  The question is not whether information was discarded.
  The question is whether the discarded information is
  \emph{relevant to the reconstruction task}.
  A road map that omits bird migration patterns is not
  deficient for a traveler.
  It is appropriately filtered.
\end{Misconception}

\principle{
  The first lesson of admissibility is not that information is
  lost.\\[4pt]
  It is that information \emph{must} be lost.\\[4pt]
  Without forgetting, representation itself would be impossible.
}

\section{Many maps, one territory}

Different projections of the same territory reveal different
structure.
No single projection is complete.

This suggests a natural question: given a collection of maps,
can the territory be recovered?

\begin{example}[Recovering terrain from multiple projections]
  A topographic map preserves elevation.
  A road map preserves connectivity.
  A geological map preserves rock type.
  Given all three simultaneously, a geologist can reconstruct
  considerably more of the territory than any single map provides.
  The reconstruction is still incomplete---the maps say nothing
  about weather, population, or history---but it is richer than
  any individual projection.
\end{example}

\begin{definition}[Reconstruction from projections]
  \label{def:reconstruction}
  Given a collection of projections
  $\{\pi_i : X \to M_i\}_{i \in I}$
  and observed data $\{m_i \in M_i\}_{i \in I}$,
  the \emph{reconstruction problem} asks:
  find $x \in X$ such that $\pi_i(x) = m_i$ for all $i \in I$.
\end{definition}

In general, the reconstruction problem has no unique solution.
The set of solutions is
\[
  \mathcal{R}(\{m_i\})
  = \bigcap_{i \in I} \pi_i^{-1}(m_i)
  \subseteq X.
\]
The reconstruction is unique if and only if $\mathcal{R}$ contains
exactly one point.

When $\mathcal{R}$ contains many points, additional constraints
are required to select among them.
These constraints are \emph{admissibility conditions}.

\begin{AdmPreview}
  The reconstruction problem above already has the structure
  that will drive the entire book.
  Given multiple incomplete observations, we ask:
  which original states are consistent with all observations
  simultaneously?
  The set of consistent states is the admissibility class.
  Choosing among them requires criteria---admissibility
  conditions---that specify which reconstructions are acceptable.
  This idea will be made mathematically precise in Part~V
  using sheaf theory and cohomology.
\end{AdmPreview}

\section{The holonomic perspective}

The title of this book takes its name from a paper by Solórzano
that formalizes precisely the situation we have just described.

A \emph{holonomic space} in Solórzano's sense is a triple
$(V, H, L)$ where $V$ is a normed vector space, $H$ is a subgroup
of norm-preserving automorphisms of $V$, and $L : H \to \mathbb{R}$
is a group-norm measuring the cost of transformations.
The associated metric is
\[
  d_L(u, v)
  = \inf_{a \in H} \sqrt{L^2(a) + \|u - av\|^2}.
\]

In plain language: the distance between two states is not just
how far apart they are, but how far apart they are after allowing
admissible internal reconfigurations at a measurable cost.

For the map problem, the holonomic interpretation is:
\begin{itemize}[leftmargin=2em]
  \item $V$ is the space of territory states,
  \item $H$ is the group of map-preserving rearrangements of the
    territory (transformations that do not change the map image),
  \item $L(h)$ measures the cost of such a rearrangement.
\end{itemize}
Two territory states that differ only by a map-preserving
rearrangement are holonomically close even if they differ in
absolute terms.

This will become technically important in
Chapter~\ref{ch:sheaf}, where each open observational domain
carries a holonomic triple rather than a bare vector space.

\begin{ReconExercise}
  Suppose two cities $x_1$ and $x_2$ are mapped to the same
  point $m \in M$ by a road map $\pi$.
  \begin{enumerate}[leftmargin=2em]
    \item What information about $x_1$ and $x_2$ can still be
      recovered from $m$ alone?
    \item What information is provably unrecoverable from $m$
      alone?
    \item If a second map $\pi'$ (e.g., a topographic map)
      assigns $x_1$ and $x_2$ to \emph{different} points,
      what does the pair $(\pi, \pi')$ allow you to reconstruct?
  \end{enumerate}
  This exercise previews the multi-projection reconstruction
  problem that will recur throughout the book.
\end{ReconExercise}

\section{What this chapter anticipates}

This chapter has introduced three ideas that will recur
throughout the book.

\begin{enumerate}[leftmargin=2em]
  \item \textbf{Projection.}
    Every representation is a map from a richer domain to a
    poorer one: $\pi : X \to M$.

  \item \textbf{Nontrivial kernel.}
    Every useful representation discards information:
    $\ker(\pi) \neq \emptyset$.

  \item \textbf{Reconstruction.}
    Given multiple projections, one can attempt to recover the
    original structure, but the recovery is generally incomplete
    or non-unique.
\end{enumerate}

\begin{MathEcho}
  This chapter introduced:
  \[
    \pi : X \to M,
    \qquad
    \ker(\pi),
    \qquad
    X / {\sim_\pi},
    \qquad
    \mathcal{R}(\{m_i\}) = \bigcap_i \pi_i^{-1}(m_i).
  \]
  These objects will reappear in:
  \begin{itemize}[leftmargin=1.5em,itemsep=0pt]
    \item Chapter~\ref{ch:detector}
      (observational equivalence classes),
    \item Chapter~\ref{ch:horizons}
      (horizon-relative reconstruction),
    \item Chapter~\ref{ch:clio}
      (reconstruction error and CLIO),
    \item Chapter~\ref{ch:sheaf}
      (admissibility sheaf and gluing),
    \item Chapter~\ref{ch:monoturn}
      (the Monoturn as global reconstruction target).
  \end{itemize}
\end{MathEcho}

\begin{OpenQuestion}
  Given only a map $\pi : X \to M$ and no direct access to
  the territory $X$, what classes of territory can be uniquely
  reconstructed from $M$?

  This question appears trivial in cartography.
  It becomes fundamental in inverse problems, detector theory,
  and cosmology, where $X$ is never directly accessible and
  $M$ is all the observer ever has.
\end{OpenQuestion}

\WhatSurvived{Road map $\pi : X \to M$}{%
  Connectivity between locations\\
  Approximate relative position\\
  Route structure\\
  Named landmarks%
}{%
  Elevation\\
  Physical texture\\
  Most geometry\\
  Temporal history\\
  Unlabeled features%
}{%
  Partial. Combining multiple map types\\
  improves reconstruction.%
}{%
  $\ker(\pi)$: all pairs $(x_1, x_2)$ with\\
  $\pi(x_1) = \pi(x_2)$, including distinct\\
  locations mapped to the same road node.%
}

\begin{HolonomyNote}
  After traversing this chapter, what information appears only
  in retrospect?

  The equivalence relation $x_1 \sim_\pi x_2$ looks like a
  simple identification.
  But following this chapter's loop---projection, kernel,
  quotient, reconstruction, holonomic metric---reveals that the
  equivalence relation does not merely erase structure.
  It \emph{defines} the class of transformations that are
  invisible to the observer.

  The holonomy of a representation is precisely the group of
  transformations that leave all observations unchanged.
  A large holonomy group means many invisible transformations
  and a severely underdetermined reconstruction problem.
  A trivial holonomy group means the representation is injective
  and reconstruction is unique.

  Every chapter in this book is a study of the holonomy of
  some observation system.
\end{HolonomyNote}
